\begin{table}[htbp]
   \caption{Where the Fall in On-Time Assignment Happens}\label{tab: gate decomposition}
   \centerline{\scalebox{0.9}{\begin{tabular}{lccc}
      \toprule
                   & \multicolumn{2}{c}{Pass rate (per 100 reaching)} & Contribution\\
                   & Winners & Losers & (seats per 100)\\
      \midrule
      Finished school on time & 78.17 & 80.29 & -0.7361$^{***}$\\
      Has a usable exam score & 73.18 & 73.01 & +0.0644\\
      Submitted an application & 61.71 & 62.00 & -0.1249\\
      Selected where she listed & 73.72 & 76.63 & -1.0249$^{***}$\\
      \midrule
      Gets a seat (all four) & 26.03 & 27.85 & -1.8215$^{***}$\\
      \bottomrule
   \end{tabular}}}
   \par\smallskip{\footnotesize Notes: the chance of a seat is the product of the four pass rates; the winner--loser difference in the product is divided into one additive piece per step by substituting one rate at a time in time order (Kitagawa 1955; Das Gupta 1993; ordering per Kim and Strobino 1984; funnel-decomposition precedent Chetty, Deming and Friedman 2026). A usable exam score requires both compulsory tests. Rates are measured within each draw and averaged with precision weights; the pieces sum to the difference exactly (clustered SE of the total 0.4278). Counted strictly step by step, 11,925 of the certified panel's 12,006 assigned children clear all four steps; the 81 remaining are recorded with a seat despite failing an earlier step on paper, which is the whole gap between this table's total and the certified assignment row ($-$1.8356 seats per 100). Stars from randomization inference (9,999 within-draw reshuffles of who wins, the whole split recomputed on each reshuffle): $^{***}$p$<$0.01, $^{**}$p$<$0.05, $^{*}$p$<$0.1.}
\end{table}

